\documentclass[conference]{IEEEtran}
\usepackage[english]{babel}
\usepackage{cite}
\ifCLASSINFOpdf
 \usepackage[pdftex]{graphicx}
 \graphicspath{{./images/}{./svg-inkscape/}}
 \DeclareGraphicsExtensions{.pdf,.jpeg,.png}
\else
\fi
\usepackage{amsmath}
\usepackage{amssymb}
\newcommand{\argmax}{\operatorname*{argmax}}
\usepackage{array}
\usepackage{url}
\usepackage{hyperref}
\hyphenation{op-tical net-works semi-conduc-tor}
\usepackage{xcolor}
\usepackage{subcaption}

\title{Does Functional Similarity Induced by Distillation Preserve Internal Representations and Explanatory Mechanisms in Autoregressive Models?\thanks{This is a condensed, page-limited version of this paper. An extended version with complete per-layer/per-model tables, per-head and hidden-state (CKA) heatmaps, additional explanation figures, and further experimental detail is available as \texttt{main.tex} in the project repository.}}

\author{
\IEEEauthorblockN{Martín Perciante}
\IEEEauthorblockA{Universidad Católica del Uruguay\\ Montevideo, Uruguay\\\href{mailto:tincho.perciante@gmail.com
}{tincho.perciante@gmail.com}}}

\begin{document}

\maketitle

\begin{abstract}
Knowledge distillation is often evaluated in terms of functional similarity between teacher and student models, typically measured through predictive performance. However, it remains unclear whether this functional similarity also implies the preservation of internal representations and explanatory mechanisms, or whether it is simply a consequence of training a smaller model on the same task. This work investigates this question in autoregressive Transformers. A TinyLlama teacher is distilled by depth into a family of $11$-layer students spanning a non-distilled baseline, a weakly-distilled variant, and two distilled variants of increasing fidelity (standard and class-restricted $KL$), isolating the effect of distillation strength from that of depth reduction alone. Internal representations are compared through layer-wise and global attention-map metrics, per-head similarity, and linear CKA on hidden states, while explanatory mechanisms are evaluated through \textbf{LIME}, \textbf{SHAP}, and \textbf{Integrated Gradients} (zero-vector baseline), moving beyond qualitative examples to aggregate agreement and faithfulness metrics. The study is conducted first as a pilot on binary IMDb sentiment classification and then, as the main experiment, on four-class AG News topic classification. Distillation does not transfer the teacher's explanatory mechanisms on either dataset. Its effect on classification accuracy is task-dependent (it helps on IMDb but not on AG News, where a non-distilled baseline outperforms every distilled variant), yet it robustly and monotonically improves distributional fidelity to the teacher on both datasets regardless of accuracy. Representational similarity is only partially, and inconsistently, attributable to distillation once a non-distilled baseline is available for comparison. These results indicate that classification performance, distributional similarity, explanatory alignment, and representational alignment are largely dissociable properties of a distilled model.
\end{abstract}

\section{Introduction}

Explainable Artificial Intelligence (XAI) seeks to interpret and explain the behavior of complex models, which is essential in application domains where deployed systems must justify autonomous decisions and foster user trust. Autoregressive Transformer-based models have driven recent progress in Natural Language Processing (NLP), but their computational cost has motivated compression techniques such as knowledge distillation (KD), in which a smaller student model is trained to approximate the output distribution of a larger teacher \cite{hinton2015distilling}.

Although KD targets functional behavior, it remains unclear whether it also induces similarity in the teacher's internal representations and explanatory mechanisms, or whether any such similarity is simply a byproduct of training a smaller model of the same architecture on the same task -- a question sharpened by prior work showing that different attention configurations can induce similar predictions \cite{jain2019attention}, challenging the interpretation of attention as a faithful explanation.

This work investigates whether functional similarity induced by logit-based KD extends to representational and explanatory similarity, and, critically, disentangles this from the effect of depth reduction alone. We adopt a generative classification setting \cite{daillm}, in which target classes correspond to single vocabulary tokens, allowing internal representations to be compared directly between teacher and student. We first conduct a pilot study on binary IMDb sentiment classification, then extend the full analysis, as the main experiment, to four-class, non-binary AG News topic classification.

In summary, this work makes the following contributions:
\begin{itemize}
    \item A family of four student variants sharing a single parameterized loss but spanning a non-distilled baseline, a weakly-distilled variant, a standard global-$KD$ variant, and a class-restricted $KD$ variant, disentangling the effect of distillation strength from that of depth reduction alone.
    \item A multi-level representational comparison framework: layer-wise and global attention analyses (cosine similarity, Frobenius distance, $JS$ divergence, spectral effective-rank), per-head attention similarity, and linear CKA on hidden states.
    \item A move from qualitative to systematic explanation comparison, evaluating \textbf{LIME}, \textbf{SHAP}, and \textbf{Integrated Gradients} with aggregate top-$k$ overlap, rank/value correlation, sign agreement, faithfulness, and stability metrics.
    \item Extension of the full analysis from an initial binary pilot (IMDb) to a four-class, non-binary main study (AG News), testing whether observed patterns generalize beyond a single task.
\end{itemize}

\section{Related Work}

Knowledge distillation (KD) transfers information from a high-capacity teacher to a smaller student by combining a cross-entropy term with a Kullback--Leibler ($KL$) divergence between output distributions, often softened by a temperature $T$ \cite{hinton2015distilling}. KD is widely used in deep networks and large language models, but its success is evaluated almost exclusively via functional metrics such as accuracy \cite{gou2021knowledge}, leaving open whether functional similarity coincides with similarity in internal representations or decision mechanisms. \textit{Generative classification} \cite{daillm} reformulates classification as next-token prediction restricted to single vocabulary tokens, preserving the model's original architecture and thereby enabling direct comparison of internal representations between teacher and student, which we adopt in this work.

Post-hoc explainability methods such as \textbf{LIME} \cite{ribeiro2016should} and \textbf{SHAP} \cite{lundberg2017unified} treat the model as a black box, while gradient-based methods such as \textbf{Integrated Gradients} \cite{sundararajan2017axiomatic} exploit differentiability; all three are commonly used to quantify token-level contributions to a prediction. Beyond post-hoc attribution, attention maps offer a potentially interpretable internal mechanism, but whether they constitute a faithful explanation remains contested: Jain and Wallace \cite{jain2019attention} show that alternative attention configurations can induce similar predictions, whereas Wiegreffe and Pinter \cite{wiegreffe2019attention} argue attention can still provide plausible explanations under appropriate conditions. \textbf{Attention Rollout} \cite{abnar2020quantifying} aggregates attention across layers, accounting for residual connections, to characterize global information flow. To the best of our knowledge, whether a student that closely reproduces a teacher's predictions also reproduces the internal representations and explanatory mechanisms probed by these tools -- and how much of any such reproduction is attributable to distillation itself rather than to depth reduction alone -- remains unexplored; this work addresses that gap.

\section{Experiments}

\subsection{Experimental Setup}\label{sec:setup}

The project used \textbf{Python}, \textbf{PyTorch}, Hugging Face's \textbf{transformers}, \textbf{lime}/\textbf{shap} for explanation, \textbf{SciPy} for the rank/value-correlation statistics used in the aggregate explanation metrics, and \textbf{scikit-learn} for evaluation metrics.

The teacher is \textbf{TinyLlama-1.1B-Chat-v1.0} \cite{zhang2024tinyllama}, a $22$-layer decoder-only Transformer \cite{vaswani2017attention}. Every student shares the teacher's embeddings, attention mechanism, and hidden dimensions, reducing only depth to $11$ randomly-initialized blocks ($2{:}1$ layer correspondence; compression rate $CR\approx1.787$, $\approx44\%$ fewer parameters, since embeddings and the output head are kept identical).

Classification is reformulated as generative classification \cite{daillm}: labels are single vocabulary tokens, inputs are inserted into a fixed template, and only the class-token logits are compared at inference. For \textbf{IMDb} \cite{maas-EtAl:2011:ACL-HLT2011} ($50000$ reviews, labels \textit{positive}/\textit{negative}), the template is \texttt{"Review: \{input\} Summary:"}; filtered to the teacher's $2048$-token context, splits are $22492$/$2500$/$24995$ (train/val/test). For \textbf{AG News} \cite{zhang2015character} ($4$ classes: \textit{World}, \textit{Sports}, \textit{Business}, \textit{Sci/Tech}), the template is \texttt{"Article: \{input\}\textbackslash n\textbackslash nTheme:"}; since \textit{Sci/Tech} tokenizes to four sub-tokens, it is relabeled \textit{Science} (one token) purely to preserve the single-token framing. Splits are $30000$/$100$/$7600$ (balanced). All other hyperparameters are shared across datasets.

The teacher is fine-tuned per dataset for one epoch (lr $2\times10^{-5}$, batch size $1$, AdamW, early stopping after $5$ non-improving validation rounds). Every student is randomly initialized and trained for $3$ epochs (lr $1\times10^{-4}$, no early stopping), minimizing

\[
\mathcal{L}
=
\alpha\,\mathcal{L}_{CE}
+
(1-\alpha)\,T^2 D_{KL}(P_P \Vert P_S)
+
\beta\,T^2 D_{KL}(P_P^{(c)} \Vert P_S^{(c)})
\]

with $T=2$ fixed, $P_P^{(c)},P_S^{(c)}$ the distributions restricted to the classification-relevant tokens, and $(\alpha,\beta)$ set per Table~\ref{tab:model_family}, defining a family of four students spanning no distillation ($B$) to full class-restricted distillation ($S_2$).

\begin{table}[h!]
\centering
\caption{Model family. All students share the same reduced ($11$-layer) architecture and differ only in $(\alpha,\beta)$.}
\small
\begin{tabular}{|l|c|c|c|l|}
\hline
Model & Init. & $\alpha$ & $\beta$ & Description \\ \hline
$P$ (teacher) & pretrained & -- & -- & fine-tuned, CE only \\
$B$ & random & $1$ & $0$ & no distillation \\
$S_{bad}$ & random & $0.9$ & $0$ & weak global $KD$ \\
$S_1$ & random & $0.5$ & $0$ & standard global $KD$ \\
$S_2$ & random & $0.5$ & $0.25$ & $+$ class-restricted $KD$ \\ \hline
\end{tabular}
\label{tab:model_family}
\end{table}

$B$ isolates how much teacher-student similarity is attributable to distillation itself versus simply training a smaller model on the same task; $S_{bad}$ provides an intermediate point. $S_1$ uses only the global $KD$ term; since global $KL$ divergence does not guarantee agreement on the $\argmax$ among target-class logits specifically, $S_2$ adds the class-restricted term, weighted lightly to act as a regularizer.

\textbf{Integrated Gradients} uses a zero embedding vector as its baseline (not the padding token used in earlier iterations of this study). Functional evaluation compares teacher and student explanations both qualitatively, on one illustrative example per dataset, and systematically, via the aggregate metrics of Section~\ref{sec:metrics} computed over $N=100$ held-out examples per dataset. Representational comparisons use attention maps averaged over held-out samples ($N=200$), individual attention heads ($N=50$ prompts), and full hidden states via linear CKA ($N=100$ prompts). Experiments ran on an NVIDIA RTX 5070 Ti (16GB). Code and the extended version of this paper: \textbf{https://github.com/PseudoKiwi/ProyectoXAI}. Manuscript writing was assisted by Claude Sonnet 4.6 \cite{anthropic2025claude} and ChatGPT \cite{openai_chatgpt}, used exclusively for stylistic revision.

\subsection{Comparison Metrics}\label{sec:metrics}

Attention maps are row-stochastic matrices; we compare them with complementary geometric, probabilistic, and spectral measures. Given two attention matrices $A,B$, \textit{cosine similarity} is $\cos(\theta)=\langle A,B\rangle/(\|A\|_2\|B\|_2)\in[-1,1]$, and the \textit{Frobenius distance} is $d_F(A,B)=\|A-B\|_F$, bounded by $\sqrt{2n}$ for row-stochastic $n\times n$ matrices ($64$ for our $n=2048$ maps). The row-averaged \textit{Jensen--Shannon divergence} is bounded in $[0,\ln 2]$ \cite{lin1991divergence}. Spectral complexity uses the entropy-based \textit{effective rank} \cite{roy2007effective} and the \textit{Participation Ratio} \cite{recanatesi2022scale}. The same formulas are additionally applied to individual attention heads rather than head-averaged maps. Hidden-state representations (all layers, all models) are compared directly via linear \textit{Centered Kernel Alignment} (CKA) \cite{kornblith2019similarity}, a rotation/scale-invariant similarity measure bounded in $[0,1]$.

To move beyond qualitative explanation comparison, we compute aggregate metrics over held-out samples, comparing each student's word-level attribution against the teacher's: top-$k$ \textit{Jaccard overlap} ($k=10$) between the top-$k$ attributed words; \textit{Spearman/Pearson correlation} of signed scores over the common attributed vocabulary; \textit{sign agreement@$k$}; deletion-based \textit{faithfulness@$k$} \cite{deyoung2020eraser}, the drop in true-class probability after masking the top-$k$ words (computed per model, including the teacher, as a reference); and \textit{stability}, the mean pairwise Jaccard@10 across $5$ re-runs with different seeds.

\subsection{Evaluation on the \textbf{Test} Set}\label{sec:test-eval}

Table~\ref{tab:perf} reports classification performance for all five models on both datasets. On \textbf{IMDb}, distillation clearly helps: $S_2$ recovers most of $S_1$'s deficit and approaches the teacher, while the non-distilled $B$ lags behind. On \textbf{AG News}, the opposite holds: $B$ achieves the \emph{highest} accuracy of all five models, ahead of every distilled variant and the teacher ($B>S_{bad}>S_2>S_1>P$) -- teacher accuracy is limited chiefly by confusing \textit{Business} with \textit{Science} ($344$ of $1900$ Business articles), a confusion every student reduces without eliminating.

\begin{table}[h!]
\centering
\caption{Accuracy / F1 (macro-averaged for AG News) on the \textbf{test} set.}
\small
\begin{tabular}{|l|l|c|c|c|c|c|}
\hline
Dataset & Metric & $P$ & $B$ & $S_{bad}$ & $S_1$ & $S_2$ \\ \hline
IMDb & Acc. & $0.968$ & $0.858$ & $0.813$ & $0.672$ & $0.934$ \\
IMDb & F1 & $0.968$ & $0.857$ & $0.810$ & $0.636$ & $0.934$ \\ \hline
AG News & Acc. & $0.849$ & $0.886$ & $0.883$ & $0.878$ & $0.880$ \\
AG News & F1 & $0.848$ & $0.886$ & $0.884$ & $0.878$ & $0.881$ \\ \hline
\end{tabular}
\label{tab:perf}
\end{table}

Table~\ref{tab:kl} reports average $KL$ divergence to the teacher at temperatures $T=1$ and $T=2$. Unlike accuracy, $KL$ fidelity follows the same monotonic pattern on \emph{both} datasets and at \emph{both} temperatures: $B$ is by far the least faithful, and fidelity increases with distillation strength ($B>S_{bad}>S_1\gtrsim S_2$) -- distillation's most robust, reproducible effect, independent of whether it improves accuracy. Notably, on IMDb $S_1$ attains \emph{lower} $KL$ than $S_2$ at both temperatures despite far worse accuracy: global $KL$ divergence is not a predictive proxy for classification performance, nor, as shown below, of explanatory or representational alignment.

\begin{table}[h!]
\centering
\caption{Average $KL$ divergence to the teacher $P$ at $T=1$ and $T=2$ (IMDb: $N=3000$ balanced subset; AG News: full \textbf{test} set).}
\small
\begin{tabular}{|l|l|c|c|c|c|}
\hline
Dataset & $T$ & $B$ & $S_{bad}$ & $S_1$ & $S_2$ \\ \hline
IMDb & $1$ & $2.847$ & $0.104$ & $0.104$ & $0.111$ \\
IMDb & $2$ & $1.686$ & $0.132$ & $0.127$ & $0.193$ \\
AG News & $1$ & $2.080$ & $0.630$ & $0.279$ & $0.248$ \\
AG News & $2$ & $1.784$ & $0.418$ & $0.229$ & $0.247$ \\ \hline
\end{tabular}
\label{tab:kl}
\end{table}

\subsection{Functional Evaluation and Explanation Comparison}\label{sec:functional-eval}

\textbf{LIME} explanations for a representative AG News \textit{Business} example (the dataset's weakest class) show the teacher's attributions concentrated on topically relevant words, while all four students assign high importance to task-irrelevant words -- the same pattern found on the IMDb pilot (teacher highlights genuine sentiment words; students favor function words like \textit{and}/\textit{the}, regardless of classification performance); see the extended version for the corresponding figure.

Table~\ref{tab:lime_agg} reports aggregate \textbf{LIME} agreement (vs.\ teacher) and faithfulness on both datasets, $N=100$, $k=10$. On both, $S_2$ is consistently the best-aligned and most-faithful student, though the gap to $B$ is modest.

\begin{table}[h!]
\centering
\caption{\textbf{LIME} Jaccard@10, Spearman $\rho$, and Faithfulness@10 vs.\ teacher, both datasets ($N=100$, $k=10$).}
\footnotesize
\begin{tabular}{|l|c|c|c|c|c|c|}
\hline
& \multicolumn{3}{c|}{IMDb} & \multicolumn{3}{c|}{AG News} \\
Model & Jac.@10 & $\rho$ & Faithf. & Jac.@10 & $\rho$ & Faithf. \\ \hline
$P$ & -- & -- & $0.172$ & -- & -- & $0.138$ \\
$B$ & $0.183$ & $0.248$ & $0.083$ & $0.336$ & $0.321$ & $0.306$ \\
$S_{bad}$ & $0.170$ & $0.216$ & $0.202$ & $0.348$ & $0.280$ & $0.402$ \\
$S_1$ & $0.136$ & $0.189$ & $0.177$ & $0.342$ & $0.307$ & $0.299$ \\
$S_2$ & $0.212$ & $0.299$ & $0.267$ & $0.358$ & $0.343$ & $0.432$ \\ \hline
\end{tabular}
\label{tab:lime_agg}
\end{table}

\textbf{SHAP} (Table~\ref{tab:shap_agg}) shows the same $S_2$-best ranking, but faithfulness near zero for \emph{every} model, including the teacher, on both datasets -- possibly a limitation of masking-based perturbation for this architecture rather than a genuine absence of faithful explanations.

\begin{table}[h!]
\centering
\caption{\textbf{SHAP} Jaccard@10, Spearman $\rho$, and Faithfulness@10 vs.\ teacher, both datasets ($N=100$, $k=10$).}
\footnotesize
\begin{tabular}{|l|c|c|c|c|c|c|}
\hline
& \multicolumn{3}{c|}{IMDb} & \multicolumn{3}{c|}{AG News} \\
Model & Jac.@10 & $\rho$ & Faithf. & Jac.@10 & $\rho$ & Faithf. \\ \hline
$P$ & -- & -- & $0.118$ & -- & -- & $0.003$ \\
$B$ & $0.117$ & $0.239$ & $0.025$ & $0.315$ & $0.270$ & $-0.003$ \\
$S_{bad}$ & $0.123$ & $0.238$ & $0.093$ & $0.317$ & $0.275$ & $0.000$ \\
$S_1$ & $0.096$ & $0.245$ & $0.058$ & $0.320$ & $0.267$ & $0.000$ \\
$S_2$ & $0.176$ & $0.333$ & $0.120$ & $0.311$ & $0.262$ & $-0.001$ \\ \hline
\end{tabular}
\label{tab:shap_agg}
\end{table}

\textbf{Integrated Gradients} (Table~\ref{tab:ig_agg}) shows high top-$10$ token overlap ($\approx0.78$--$0.89$ across all models and both datasets) alongside near-zero or negative rank correlations, indicating agreement on \emph{which} tokens matter without agreement on their relative importance -- unlike LIME/SHAP, no student is consistently best-aligned. Sign agreement@$10$ (not tabulated for space) mirrors this split: for LIME/SHAP it follows the same $S_2$-best pattern, ranging $0.72$--$0.82$ across models and datasets, whereas for IG it sits at chance level ($0.50$--$0.54$) for every model -- IG's high token-overlap is therefore not accompanied by agreement on \emph{either} relative importance or polarity, reinforcing that its apparent cross-model consistency is comparatively shallow.

\begin{table}[h!]
\centering
\caption{\textbf{Integrated Gradients} (zero baseline) Jaccard@10, Spearman $\rho$, and Faithfulness@10 vs.\ teacher, both datasets ($N=100$, $k=10$).}
\footnotesize
\begin{tabular}{|l|c|c|c|c|c|c|}
\hline
& \multicolumn{3}{c|}{IMDb} & \multicolumn{3}{c|}{AG News} \\
Model & Jac.@10 & $\rho$ & Faithf. & Jac.@10 & $\rho$ & Faithf. \\ \hline
$P$ & -- & -- & $0.010$ & -- & -- & $0.054$ \\
$B$ & $0.793$ & $-0.005$ & $0.010$ & $0.891$ & $0.022$ & $0.067$ \\
$S_{bad}$ & $0.778$ & $0.028$ & $0.085$ & $0.883$ & $0.040$ & $0.163$ \\
$S_1$ & $0.814$ & $0.010$ & $0.025$ & $0.876$ & $-0.036$ & $0.138$ \\
$S_2$ & $0.787$ & $0.039$ & $0.019$ & $0.888$ & $0.013$ & $0.340$ \\ \hline
\end{tabular}
\label{tab:ig_agg}
\end{table}

Table~\ref{tab:stability} reports \textbf{LIME} stability (mean pairwise Jaccard@10 across $5$ re-runs, $N=20$). The teacher is more stable than the students on AG News but not always on IMDb, so stability alone is not a proxy for explanation quality.

\begin{table}[h!]
\centering
\caption{\textbf{LIME} stability (mean pairwise Jaccard@10 across $5$ seeds, $N=20$), both datasets.}
\small
\begin{tabular}{|l|c|c|c|c|c|}
\hline
Dataset & $P$ & $B$ & $S_{bad}$ & $S_1$ & $S_2$ \\ \hline
IMDb & $0.548$ & $0.587$ & $0.684$ & $0.611$ & $0.710$ \\
AG News & $0.774$ & $0.669$ & $0.639$ & $0.615$ & $0.688$ \\ \hline
\end{tabular}
\label{tab:stability}
\end{table}

\subsection{Representational Comparison}\label{sec:repr-comparison}

Table~\ref{tab:attn} reports the cosine similarity between each student's attention map and the teacher's average attention map ($P_{avg}$, matched layer pairs), at three representative layers per dataset (full per-layer tables are in the extended version). On \textbf{IMDb}, $B$ is consistently \emph{least} similar to the teacher from layer $1$ onward ($0.13$--$0.29$ across all $11$ layers, vs.\ $0.19$--$0.69$ for the three distilled variants), suggesting distillation does confer attention-level alignment beyond depth reduction alone on this task. On \textbf{AG News}, this does \emph{not} replicate: $B$ is comparable to, and at several layers (e.g.\ layer $4$: $B=0.52$ vs.\ $0.29$--$0.35$ for the distilled students) more similar to the teacher than any distilled variant -- the sharpest finding of this work regarding representational transfer.

\begin{table}[h!]
\centering
\caption{Cosine similarity between each student's attention map and the teacher's $P_{avg}$, at representative layers $i$.}
\small
\begin{tabular}{|l|c|c|c|c|c|}
\hline
Dataset & $i$ & $B$ & $S_{bad}$ & $S_1$ & $S_2$ \\ \hline
IMDb & $0$ & $0.933$ & $0.923$ & $0.937$ & $0.923$ \\
IMDb & $5$ & $0.182$ & $0.194$ & $0.394$ & $0.303$ \\
IMDb & $10$ & $0.135$ & $0.271$ & $0.234$ & $0.164$ \\
AG News & $0$ & $0.981$ & $0.965$ & $0.968$ & $0.975$ \\
AG News & $4$ & $0.518$ & $0.345$ & $0.290$ & $0.345$ \\
AG News & $10$ & $0.311$ & $0.278$ & $0.226$ & $0.264$ \\ \hline
\end{tabular}
\label{tab:attn}
\end{table}

On both datasets, Frobenius distance and $JS$ divergence are nearly indistinguishable across all four models at every layer (e.g.\ IMDb Frobenius at $i=0$: $1.24$--$1.36$ across all models, rising to $34$--$39$ around $i=2$--$3$, out of a theoretical maximum of $64$), and global (whole-model) comparisons show the same pattern as the layer-wise ones, with $B$ never clearly worse than the distilled students (e.g.\ AG News, rollout vs.\ $P_{avg}$: $B=0.986$, $S_1=0.979$, $S_2=0.975$). At an even finer grain, per-head attention-similarity heatmaps and full hidden-state linear-CKA heatmaps (Section~\ref{sec:metrics}) show $B$ visually indistinguishable from the three distilled students on \emph{both} datasets -- the clearest evidence that hidden-state and per-head geometric similarity to the teacher is largely insensitive to whether a model was distilled at all (see the extended version for the corresponding heatmap figures).

Tables~\ref{tab:pr} and~\ref{tab:er} report the spectral effective-rank measures on IMDb at the same representative layers (AG News shows the same qualitative pattern -- students intermediate between a near-rank-one $P_{avg}$ and a highly dispersed $P_r$ ($ER\approx1650$--$2026$), with no consistent ordering by distillation strength -- but is omitted here for space). $P_{avg}$ stays close to rank-one throughout ($PR\approx1$ beyond $i=0$), while $P_r$ is far more dispersed; all four students sit in between, with effective ranks generally growing toward deeper layers and no model consistently highest or lowest -- e.g.\ $B$ has the lowest $ER$ at $i=0$ ($60.2$) but is mid-pack by $i=10$. As with cosine similarity, spectral complexity is therefore not obviously ordered by distillation strength either.

\begin{table}[h!]
\centering
\caption{Participation Ratio ($PR$) on IMDb, at representative layers $i$.}
\footnotesize
\begin{tabular}{|c|c|c|c|c|c|c|}
\hline
$i$ & $B$ & $S_{bad}$ & $S_1$ & $S_2$ & $P_{avg}$ & $P_r$ \\ \hline
$0$ & $2.6$ & $2.6$ & $2.5$ & $2.7$ & $3.3$ & $399.6$ \\
$5$ & $5.6$ & $17.2$ & $5.9$ & $7.7$ & $1.0$ & $1.7$ \\
$10$ & $9.7$ & $8.6$ & $17.0$ & $12.5$ & $1.0$ & $1.9$ \\ \hline
\end{tabular}
\label{tab:pr}
\end{table}

\begin{table}[h!]
\centering
\caption{Entropy-based effective rank ($ER$) on IMDb, at representative layers $i$.}
\footnotesize
\begin{tabular}{|c|c|c|c|c|c|c|}
\hline
$i$ & $B$ & $S_{bad}$ & $S_1$ & $S_2$ & $P_{avg}$ & $P_r$ \\ \hline
$0$ & $60$ & $117$ & $96$ & $107$ & $309$ & $2027$ \\
$5$ & $356$ & $507$ & $470$ & $489$ & $26$ & $1798$ \\
$10$ & $491$ & $509$ & $551$ & $528$ & $129$ & $1819$ \\ \hline
\end{tabular}
\label{tab:er}
\end{table}

\subsection{Discussion}\label{sec:discussion}

Taken together, Sections~\ref{sec:test-eval}--\ref{sec:repr-comparison} suggest a practical reframing of the question motivating the baseline models introduced in this revision. If the only goal is task accuracy under a fixed, reduced architecture, training from scratch with cross-entropy alone ($B$) is a strong, cheaper baseline that can match or beat distillation depending on the task -- AG News being the clearest counter-example to the intuition that distillation should always help. Distillation remains preferable when the goal is matching the teacher's \emph{output distribution} (e.g.\ soft-label pipelines, calibration transfer, ensembling), where its advantage over $B$ is large and consistent (Table~\ref{tab:kl}), and it offers a modest but consistent edge in \textbf{LIME}/\textbf{SHAP} alignment even where it does not help accuracy.

More broadly, our results caution against treating distributional or representational similarity to the teacher as a proxy for task performance: on IMDb, $S_1$ attains \emph{lower} $KL$ than $S_2$ while performing far worse (Tables~\ref{tab:perf}--\ref{tab:kl}); on AG News, $B$ attains the highest accuracy while showing no advantage -- and at several layers a disadvantage -- in attention-level similarity to the teacher (Table~\ref{tab:attn}). Neither similarity reliably tracks performance here, echoing the point raised for explanations in Section~\ref{sec:functional-eval}: $S_2$ is the best-performing student on IMDb and the best-aligned on \textbf{LIME}/\textbf{SHAP}, yet trails $S_1$ in $KL$ fidelity.

\subsection{Does Distillation Add Value Beyond Depth Reduction?}\label{sec:distillation-value}

Table~\ref{tab:verdict} summarizes the picture across both datasets. Distillation's only robust, task-independent effect is pulling the student's output distribution toward the teacher's; its effects on accuracy, explanatory alignment, and fine-grained representational similarity are comparatively small, inconsistent across datasets, and, for attention/hidden-state geometry, largely absent once a non-distilled baseline is available for comparison.

\begin{table}[h!]
\centering
\caption{Does distillation help beyond training a smaller model?}
\small
\begin{tabular}{|l|l|l|}
\hline
Property & IMDb & AG News \\ \hline
Accuracy & Helps ($S_2\gg B$) & Hurts ($B>$ all) \\
$KL$ fidelity & Robustly helps & Robustly helps \\
Explanation align. & Modest ($S_2$ best) & Modest ($S_2$ best) \\
Represent. align. & Partial ($B$ worst) & Absent ($B$ competitive) \\ \hline
\end{tabular}
\label{tab:verdict}
\end{table}

\subsection{Limitations}\label{sec:limitations}

This study covers only two English text-classification datasets of modest label cardinality (binary IMDb, four-class AG News); whether the task-dependence between them (Table~\ref{tab:verdict}) generalizes to other modalities or higher-cardinality tasks is untested. The $(\alpha,\beta)$ grid is narrow: only two points beyond $S_1$/$S_2$ were tested ($\alpha=1$ for $B$, $\alpha=0.9$ for $S_{bad}$), and the opposite extreme -- pure distillation with no cross-entropy term ($\alpha=0$) -- was never tested (Section~\ref{sec:distillation-value}). The near-zero \textbf{SHAP} faithfulness and \textbf{IG} high-overlap/low-rank-correlation pattern, observed for every model including the teacher, may partly reflect limitations of these methods rather than of distillation. Underlying all of this is a hardware constraint: every experiment ran on a single consumer GPU (RTX 5070 Ti, $16$GB), bounding how many datasets, scales, and teacher/student architectures could be explored in reasonable time -- likely a more binding constraint than the use of a single teacher architecture (TinyLlama-1.1B) throughout.

\section{Conclusions}

This work investigated whether functional similarity induced by knowledge distillation extends to internal representations and explanatory mechanisms, using a family of models -- a non-distilled baseline, a weakly-distilled variant, and two distilled students -- on an IMDb pilot and, as the main study, on AG News. In these specific, comparatively simple settings, the answer is negative rather than definitively so: every student, including the non-distilled baseline, assigns importance to irrelevant tokens regardless of accuracy, and $B$ is indistinguishable from the distilled students at the level of individual attention heads and hidden-state geometry on both datasets, and even at the level of full attention maps, distillation helps only on IMDb. The one property that behaves consistently, independently of accuracy, is distributional ($KL$) fidelity to the teacher.

A striking pattern underlies this: $B$, $S_{bad}$, $S_1$, and $S_2$ converge toward similar attention and hidden-state structure despite spanning $\alpha$ from $0.5$ to $1$. Since all four are trained from scratch on the same architecture and dataset, and the one term every one of them shares is cross-entropy, this convergence is at least as consistent with the task and dataset imprinting this structure as with a distillation-specific effect; our design cannot distinguish the two, since it never removes cross-entropy from the loss entirely. A natural next step, left to future work, is a student trained with \emph{pure} distillation ($\alpha=0$, no cross-entropy term): if its representations diverge from the $B$/$S_{bad}$/$S_1$/$S_2$ cluster, that would implicate cross-entropy, rather than the dataset alone, as the source of this convergence.

\bibliographystyle{plain}
\bibliography{references}

\end{document}
